\subsection{Joint spectra of stochastic matrices: a bounded development of GP-1987-21}
\label{dev:stochastic}

\paragraph{The source question and its scope.}
In item \hyperref[item:1987-728.4]{1987-728.4}, Gelfand asks for the
\ru{набор соб. значений в \(\mathbb C^n\)} and emphasizes that projection onto
one coordinate may conceal its geometry. The scan of P001418 supports a
question about all eigenvalues of a real nonnegative row-stochastic matrix.
It does not impose symmetry, double stochasticity, irreducibility, a specified
zero pattern, or a prescribed Jordan form. Those are different problems.
The strongest precise formulation supported by the page is to describe,
for each positive integer \(n\), the permutation-invariant set of lists
for which there exists a matrix \(P\) satisfying
\[
 \Sigma_n=\left\{(\lambda_1,\ldots,\lambda_n)\in\mathbb C^n:
 \begin{array}{l}
 \det(tI-P)=\prod_{j=1}^n(t-\lambda_j),\\
 P\in\mathbb R^{n\times n},\quad
 P_{ij}\geq0,\quad P\mathbf1=\mathbf1
 \end{array}\right\}.
\]
Multiplicities are retained. This is a reconstruction of a shorthand
question, not a conjectured theorem attributed to the speaker.

\paragraph{State of the art and the bounded contribution.}
The order-three characterization predates the seminar:
\citet{loewylondon1978inverse} appeared in 1978; its attribution and
explicit alternatives are recorded in \citet[p.~506]{ranteng2024patterns}.
The latter paper additionally classifies prescribed order-three zero patterns
and appeared in July 2024 (received May 2023).
\citet[Theorems~3.9--3.12 and Sections~4--5, 8]{johnsonpaparella2025perron}
establish geometric properties of joint stochastic spectra and polyhedral
families from fixed diagonalizing matrices; their preprint first circulated
on 12 September 2024 and the journal paper appeared in 2025.
Their nearly complete geometric picture in order four does not solve
all orders. As of 11 November 2026, the sources checked provide no general
structural solution in orders at least five. The distinction between
diagonalizable and unrestricted realization is substantive;
\citet[Section~2 and Proposition~1]{croninlaffey2018diagonalizable}
first circulated in January 2017 and appeared in July 2018.

For a fixed \(n\), polynomial feasibility already makes membership
decidable for algebraic input. Thus \emph{open} does not mean that there is
no finite decision procedure. What remains is an explicit, usable
description of the joint region, its boundary and realizing matrices.
The contribution here is a self-contained, checked low-dimensional atlas
and precise obstructions to extrapolating necessary tests.
None of the results below is claimed to be new.

\begin{proposition}[Geometry and a finite formulation \ours]
\label{prop:stochastic-geometry}
For every fixed \(n\), \(\Sigma_n\) is compact and semialgebraic in
\(\mathbb R^{2n}\). It is invariant under coordinate permutations and
complex conjugation, and star-shaped at \((1,\ldots,1)\).
The set
\[
 \widehat\Sigma_n=\{(\lambda_2,\ldots,\lambda_n):
                         (1,\lambda_2,\ldots,\lambda_n)\in\Sigma_n\}
\]
is star-shaped at the origin.
\end{proposition}
\begin{proof}
The stochastic matrices form a compact product of \(n\) probability
simplices. Every eigenvalue has modulus at most one: if \(Pv=\lambda v\)
and \(|v_i|=\max_j|v_j|>0\), then
\(|\lambda v_i|\leq\sum_jP_{ij}|v_j|\leq|v_i|\).
For a convergent sequence of lists in \(\Sigma_n\), choose corresponding
matrices and a convergent subsequence of these matrices. Continuity of
characteristic-polynomial coefficients shows that the limiting matrix
realizes the limiting list. Thus \(\Sigma_n\) is closed and bounded,
without requiring a continuously chosen ordering of eigenvalues.

Writing \(\lambda_j=x_j+iy_j\), equality of coefficients in
\(\det(tI-P)=\prod_j(t-x_j-iy_j)\) is a finite system of real
polynomial equalities. Append \(P_{ij}\geq0\) and
\(\sum_jP_{ij}=1\). Projecting onto the \(x,y\) coordinates and applying
real quantifier elimination gives semialgebraicity and a decision
algorithm for effectively encoded algebraic input.
The background theorem is the standard one treated in
\citet[Chapters~2, 13--14]{basupollackroy2006algorithms};
we do not reprove that theorem or claim an efficient complexity bound.

Permutation invariance follows from the characteristic polynomial;
conjugation invariance follows from its real coefficients.
For \(0\leq s\leq1\), \((1-s)I+sP\) is stochastic with eigenvalues
\(1-s+s\lambda_j\), proving the first star-shapedness assertion.
Put \(J=\mathbf1\mathbf1^{\mathsf T}/n\) and \(Q=sP+(1-s)J\).
On the invariant line spanned by \(\mathbf1\), \(Q\) is the identity.
On the quotient by that line, it induces \(s\) times the map induced
by \(P\), because \(J\) has image in the line.
In a basis beginning with \(\mathbf1\), its characteristic polynomial
is therefore \((t-1)\prod_{j=2}^n(t-s\lambda_j)\).
This also covers repeated eigenvalues and nondiagonalizable matrices.
\end{proof}

\begin{lemma}[Power-sum inequalities \ours]
\label{lem:jll}
For a list in \(\Sigma_n\), put \(s_k=\sum_j\lambda_j^k\).
Then
\[
 s_k\geq0,\qquad s_k^m\leq n^{m-1}s_{km}
 \quad(k,m\in\mathbb N).
\]
The list contains \(1\) and is closed under conjugation.
\end{lemma}
\begin{proof}
The last assertions follow from reality and \(P\mathbf1=\mathbf1\).
Let \(B=P^k\geq0\). Then \(s_k=\operatorname{tr}B\geq0\).
The term in \((B^m)_{ii}\) that remains at \(i\) at every step is
\(B_{ii}^m\); every other term is nonnegative. Consequently
\(\sum_iB_{ii}^m\leq\operatorname{tr}(B^m)=s_{km}\).
The power-mean inequality gives
\((\sum_iB_{ii})^m\leq n^{m-1}\sum_iB_{ii}^m\).
Combining the inequalities proves the claim, including \(m=1\).
\end{proof}
These are the classical Johnson--Loewy--London inequalities, independently
derived here; the historical references are discussed in
\citet[Section~1.1]{croninlaffey2018diagonalizable}.

\begin{theorem}[The complete two- and three-state regions \ours]
\label{thm:lowdim-stochastic}
A two-element multiset is the spectrum of a real row-stochastic
\(2\times2\) matrix exactly when it is \(\{1,r\}\), \(-1\leq r\leq1\).
A three-element multiset is the spectrum of a real row-stochastic
\(3\times3\) matrix exactly when it has one of the forms
\[
\begin{array}{ll}
\{1,r,s\},&r,s\in[-1,1],\quad 1+r+s\geq0;\\[2pt]
\{1,a+ib,a-ib\},& b\ne0,\quad -\tfrac12\leq a\leq1,\quad
 |b|\leq(1-a)/\sqrt3 .
\end{array}
\]
Weak inequalities retain every boundary case. The second line describes
the triangle with vertices \(1,-1/2\pm i\sqrt3/2\), with its intersection
with the real axis omitted because \(b\ne0\).
\end{theorem}
\begin{proof}
In dimension two, reality and the eigenvalue \(1\) force the other
eigenvalue to be real, and the modulus bound gives necessity.
The matrix
\[
 C(r)=\frac12\begin{pmatrix}1+r&1-r\\1-r&1+r\end{pmatrix}
\]
is nonnegative and row-stochastic for \(-1\leq r\leq1\).
Its eigenvectors \((1,1)\), \((1,-1)\) have eigenvalues \(1,r\).

In dimension three, the remaining pair is either real or conjugate.
For a real pair, the modulus and trace conditions are necessary.
For a conjugate pair, the trace gives \(a\geq-1/2\), and the
modulus bound gives \(a\leq1\).
Lemma~\ref{lem:jll} with \(k=1,m=2,n=3\) gives
\[
 (1+2a)^2\leq3\{1+2(a^2-b^2)\}
 \quad\Longleftrightarrow\quad 3b^2\leq(1-a)^2 .
\]
These are the claimed necessary inequalities.

For real-pair sufficiency, if one member is nonnegative, rename it
\(r\in[0,1]\) and write the other as \(s\in[-1,1]\). Then
\[
 P=\begin{pmatrix}
 (1+s)/2&(1-s)/2&0\\
 (1-s)/2&(1+s)/2&0\\
 1-r&0&r
 \end{pmatrix}
\]
is nonnegative and stochastic; block triangularity gives eigenvalues
\(1,s,r\). If both are negative, write \(r=-u,s=-v\), where
\(u,v>0\) and \(u+v\leq1\). The matrix
\[
 P=\begin{pmatrix}
 0&1&0\\0&0&1\\uv&u+v-uv&1-u-v
 \end{pmatrix}
\]
is stochastic and nonnegative, since
\(u+v-uv=u(1-v)+v\geq0\). Its characteristic polynomial is
\[
 t^3-(1-u-v)t^2-(u+v-uv)t-uv=(t-1)(t+u)(t+v).
\]
Zero eigenvalues were covered by the first construction.

For a conjugate pair, set
\[
 p_0=\frac{1+2a}{3},\qquad
 p_1=\frac{1-a}{3}+\frac{b}{\sqrt3},\qquad
 p_2=\frac{1-a}{3}-\frac{b}{\sqrt3}.
\]
The inequalities make these nonnegative with sum one.
For \(\omega=e^{2\pi i/3}\), the circulant
\[
 P=\begin{pmatrix}p_0&p_1&p_2\\p_2&p_0&p_1\\p_1&p_2&p_0\end{pmatrix}
\]
has eigenvectors \((1,1,1)\), \((1,\omega,\omega^2)\),
\((1,\omega^2,\omega)\), with eigenvalues
\(1,p_0+p_1\omega+p_2\omega^2,p_0+p_1\omega^2+p_2\omega\).
These are \(1,a+ib,a-ib\), completing the sufficiency proof.
\end{proof}

\begin{example}[Two insufficient tests \ours]
Each of \(-3/4\) and \(1\) occurs in a stochastic three-state
spectrum: use \(C(-3/4)\oplus(1)\).
Yet \(\{1,-3/4,-3/4\}\) cannot be a joint spectrum because its
trace is \(-1/2\). A product of marginal regions is therefore too large.
Even the modulus, conjugation and all nonnegative power-sum tests
can be insufficient. The list \(\{1,i/\sqrt2,-i/\sqrt2\}\) has
\(s_{2m+1}=1\), \(s_{2m}=1+2(-1/2)^m\geq0\), but
\(s_1^2=1>3s_2=0\). Both exclusions are exact.
\end{example}

\begin{proposition}[An order-five obstruction surviving all JLL tests \ours]
\label{prop:five-obstruction}
The list \(\Lambda=\{1,1,-2/3,-2/3,-2/3\}\) satisfies every
power-sum and JLL inequality for \(n=5\), together with the Perron and
conjugation conditions. It is nevertheless not the spectrum of any
nonnegative \(5\times5\) matrix, and hence is not in \(\Sigma_5\).
\end{proposition}
\begin{proof}
Put \(q=2/3\), so \(s_k=2+3(-q)^k\).
For even \(k\), apply the power-mean inequality to
\(1,1,q^k,q^k,q^k\), obtaining
\(s_k^m\leq5^{m-1}s_{km}\).
For odd \(k\), \(0\leq s_k\leq2\), since \(q^k\leq2/3\).
If \(m\) is odd, \(s_{km}\geq s_k\), so
\(s_k^m\leq2^{m-1}s_k\leq5^{m-1}s_{km}\).
If \(m\) is even, \(s_{km}\geq2\), so
\(s_k^m\leq2^m\leq5^{m-1}s_{km}\).
This proves all inequalities for every positive \(k,m\).

Suppose a nonnegative matrix realizes the list. Put it in Frobenius
block-triangular form, with irreducible nonnegative diagonal blocks,
including one-by-one blocks. Each block has its nonnegative spectral
radius as an eigenvalue. The only nonnegative member of \(\Lambda\)
is \(1\), so each block must contain a copy of \(1\).
The Perron eigenvalue of an irreducible block is simple; there must
therefore be exactly two blocks, each using one copy.
One block receives at least two of the three negative eigenvalues.
Its trace is at most \(1-4/3=-1/3\), impossible for a nonnegative block.
\end{proof}
This is the classical list recorded in
\citet[Section~2]{croninlaffey2018diagonalizable}.
The all-\(k,m\) check and block proof are supplied to make the
obstruction independently verifiable, without claiming novelty.

\paragraph{Independent verification \oursComputed.}
The script \path{code/verify_development.py} checks the characteristic
polynomials symbolically and then checks the constructions using a
separate exact three-by-three determinant formula on rational grids.
It also tests independently generated stochastic matrices against the
inequalities and checks the order-five example with exact fractions
over a finite exponent range. The finite tests are diagnostic;
the universally quantified claims are proved above.
The output is \path{logs/development_checks.json}.
An independent derivation of the triangular bound uses
\(\operatorname{tr}P^2=\sum_iP_{ii}^2+
2\sum_{i<j}P_{ij}P_{ji}\), followed by Cauchy--Schwarz.
The Fourier eigenvector proof checks the converse separately from
the characteristic-polynomial calculation. Both checks agree,
including boundary cases and reducible real-pair realizations.

\paragraph{What remains and what another run could accomplish.}
The bottleneck is joint compatibility: one nonnegative matrix must
realize all roots. Reducibility and repeated Perron roots constrain
how roots can be grouped. Projection to one eigenvalue, all power-sum
inequalities, or one fixed diagonalizing basis loses information.
The calculation does not resolve the general \(n\geq5\) region.
There is no evidence that this question was abandoned: the checked
2024--2025 papers show continued work.
A bounded next run could construct an exact atlas for one specified
family of four-state zero patterns, compare it with the known
order-four characterization and the constructions in
\citet[Section~8]{johnsonpaparella2025perron}, and certify every
boundary piece by a realizing matrix or polynomial obstruction.
That would provide a reusable companion/website contribution.
Its novelty would have to be established by that comparison.

\subsection{A Poisson-pencil repair for the Neumann calculation}
\label{dev:poisson-repair}

This annotation of item \hyperref[item:1987-716.6]{1987-716.6}
records a failure found by the independent check. With
\(\varepsilon_{123}=1\), the second printed table on P001391 has
\(\{M_i,M_j\}=\varepsilon_{ijk}a_kM_k\),
\(\{M_i,p_j\}=\varepsilon_{ijk}a_kp_k\), and
\(\{p_i,p_j\}=\varepsilon_{ijk}M_k\).
Its cyclic Jacobi expression at \((M_1,p_1,p_2)\) is
\((a_3-a_2)M_2\), so it is not a Poisson bracket for general
diagonal \(A\). Reversing time cannot repair this defect.
The scan confirms the printed coefficients. The following is a
consistent algebraic emendation, not a claim to recover a uniquely
intended historical formula.

\begin{proposition}[A compatible repaired bracket \ours]
For arbitrary real \(a_1,a_2,a_3\), the table
\[
 \{M_i,M_j\}_2=\varepsilon_{ijk}a_kM_k,\qquad
 \{M_i,p_j\}_2=\varepsilon_{ijk}a_jp_k,\qquad
 \{p_i,p_j\}_2=-\varepsilon_{ijk}M_k
\]
defines a linear Poisson bracket. It is compatible with the standard
\(\mathfrak e(3)^*\) bracket
\(\{M_i,M_j\}_1=\varepsilon_{ijk}M_k\),
\(\{M_i,p_j\}_1=\varepsilon_{ijk}p_k\),
\(\{p_i,p_j\}_1=0\).
With \(F=p^2/2\), the repaired bracket gives
\(\dot p=M\times p\), \(\dot M=(Ap)\times p\), exactly the
first bracket's vector field for
\(H=(M^2+p^{\mathsf T}Ap)/2\).
\end{proposition}
\begin{proof}
Let \(E_{ij}\) denote the skew matrix with entry \(1\) in position
\((i,j)\), entry \(-1\) in position \((j,i)\), and all other entries
zero. For any diagonal \(4\times4\) matrix \(D\), the product
\(X*Y=XDY\) is associative. Its commutator
\([X,Y]_D=XDY-YDX\) is consequently a Lie bracket.
Skew matrices are closed under this bracket since \(D^{\mathsf T}=D\).
Choose the basis
\[
 m_1=-E_{23},\quad m_2=-E_{31},\quad m_3=-E_{12},\qquad
 q_i=E_{i4}\quad(1\leq i\leq3).
\]
For \(D=\operatorname{diag}(a_1,a_2,a_3,-1)\), direct matrix
multiplication gives precisely the three displayed structure-constant
formulas. The associated Lie--Poisson bracket on the dual therefore
satisfies Jacobi, including when \(D\) is singular.
For \(D_0=\operatorname{diag}(1,1,1,0)\), the same basis gives the
first bracket. Every linear combination is again the commutator for
the corresponding linear combination of \(D,D_0\), proving compatibility.
Finally, contracting the repaired table with
\(\partial F/\partial p_j=p_j\) gives the two stated vector fields.
Contracting the first table with the derivatives of \(H\) gives the
same fields, independently checking their signs.
\end{proof}

\paragraph{Second check \oursComputed.}
The coordinate computation in the early-record verification script
tests every coordinate Jacobi identity for the repaired bracket and
the full pencil, and checks equality of the Hamiltonian fields.
The associative-matrix construction above is an independent proof
of those identities. Its agreement with the symbolic output confirms
the emendation; the failed original table is retained in that output.

\subsection{The complete attainable range in the tensor-angle exercise}
\label{dev:tensor-angles}

This is an elementary annotation of item
\hyperref[item:1987-737.11]{1987-737.11} and GP-1987-29, classified E;
no priority claim is made. The preamble on P001435 explicitly asks
which values the invariants in all the exercises can take. On P001436,
the tensor exercise specifies a Euclidean space \(L\) and compares
\(a\otimes L\) with \(L\otimes b\) in \(L^{\otimes3}\).
We reconstruct the finite-dimensional setting of that part of the
sheet, require \(a,b\ne0\), and normalize them to unit length;
scaling does not change either subspace. The same proof applies to
complex Hilbert spaces, although the source says Euclidean.

\begin{theorem}[Principal cosines for two tensor subspaces \ours]
Let \(L\) be a real Euclidean or complex Hilbert space of dimension
\(d\geq2\). As unit tensors \(a,b\in L\otimes L\) vary, the ordered
principal cosines \(s_1\geq\cdots\geq s_d\geq0\) between
\(a\otimes L\) and \(L\otimes b\) are exactly the lists satisfying
\[
                  \sum_{j=1}^d s_j\leq1.
\]
For \(d=1\), the sole possible list is \((1)\).
\end{theorem}
\begin{proof}
In an orthonormal basis write
\(a=\sum_{i,j}A_{ij}e_i\otimes e_j\) and
\(b=\sum_{j,k}B_{jk}e_j\otimes e_k\), so
\(\|A\|_F=\|B\|_F=1\).
The maps \(Uu=a\otimes u\), \(Vv=v\otimes b\) are isometries.
With the complex inner product conjugate-linear in its first entry,
\[
 (U^*V)_{ki}=\langle a\otimes e_k,e_i\otimes b\rangle
            =\sum_j\overline{A_{ij}}B_{jk}.
\]
Thus the cross-Gram matrix is \((\overline A B)^{\mathsf T}\),
and its singular values are the principal cosines. The trace norm
inequality
\(\|\overline A B\|_*\leq\|A\|_F\|B\|_F=1\)
gives necessity. To check this inequality directly, choose a unitary
matrix from a singular value decomposition so that its trace pairing
with \(\overline A B\) is the sum of the singular values, and apply
Cauchy--Schwarz to the matrix entries. Left or right multiplication
by a unitary matrix preserves the Frobenius norm.

Conversely, put \(T=\sum_js_j\leq1\), and let \(r\) be the number
of positive entries. All the following matrices are real, so their
constructions work over either field. If \(T=0\), take
\(A_{11}=B_{22}=1\), with all other entries zero;
both matrices have norm one and \(AB=0\).
If \(0<r<d\), set
\[
 A_{ii}=\sqrt{s_i/T},\qquad B_{ii}=\sqrt{Ts_i}
 \quad(1\leq i\leq r),\qquad
 B_{r+1,1}=\sqrt{1-T^2},
\]
and set all other entries to zero. The additional entry lies in a
row annihilated by \(A\), so
\(AB=\operatorname{diag}(s_1,\ldots,s_d)\).
Moreover, \(\|A\|_F^2=1\) and
\(\|B\|_F^2=T^2+(1-T^2)=1\).

It remains to consider \(r=d\). On the open probability simplex
\(x_j>0\), \(\sum_jx_j=1\), define
\(f(x)=\sum_js_j^2/x_j\). Cauchy--Schwarz gives \(f(x)\geq T^2\),
with equality at \(x_j=s_j/T\).
If \(T=1\), use this minimizing point. If \(T<1\), the minimum
is below one, whereas \(f(x)\) tends to infinity along a segment
approaching a face, since every \(s_j>0\).
Continuity supplies a point with \(f(x)=1\).
Now take \(A_{jj}=\sqrt{x_j}\) and
\(B_{jj}=s_j/\sqrt{x_j}\). Both Frobenius norms equal one, and
\(AB=\operatorname{diag}(s_j)\), as required.
When \(d=1\), both tensor subspaces are the entire one-dimensional
ambient space; their principal cosine is one.
\end{proof}

\paragraph{Independent verification and scope \oursComputed.}
The cross-Gram formula was re-derived entry by entry and checked
against explicit tensor embeddings, including complex coefficients.
An independent mathematical review checked all rank cases and the
simplex-continuity argument. The script
\path{code/verify_development.py} tests both real and complex
random tensors, rank-deficient constructions, and the exceptional
case \(d=1\); its output records the observed residuals.
This resolves the stated finite-dimensional attainable-range exercise.
It does not classify simultaneous invariants of several unrelated
subspaces or the Hilbert-space cases on the next part of the sheet.

\subsection{Homology spheres and twisted cohomology: checking a literal error}
\label{dev:twisted-h1}

This is an annotation of P001475, not an extra class-C or class-D
research problem. The source passes from the Hessian at a flat
\(\mathrm{SU}(2)\) connection to all local systems and asserts an
equivalence with integral first-homology vanishing.
At a specified irreducible flat connection with holonomy \(\rho\),
the relevant infinitesimal space is \(H^1(M;\operatorname{ad}\rho)\).
Ordinary and twisted homology cannot be interchanged.
Already \citet[p.~287]{atiyah1988invariants} requires perturbations
to obtain simple zeros in the homology-sphere construction.
The next two checks disprove both directions of the literal
equivalence over characteristic zero.

\begin{lemma}[Averaging one-cocycles \ours]
For a finite group \(G\) acting linearly on a real or complex vector
space \(V\), every one-cocycle \(c:G\to V\) is principal.
\end{lemma}
\begin{proof}
The identity is \(c(gh)=c(g)+g c(h)\).
Put \(S=\sum_{h\in G}c(h)\); summing gives
\(S=|G|c(g)+gS\). Thus \(c(g)=w-gw\) for \(w=S/|G|\).
The characteristic-zero hypothesis permits division by \(|G|\).
\end{proof}

\begin{example}[Twisted vanishing misses integral torsion \ours]
The closed connected orientable manifold \(M=\mathbb{RP}^3=L(2,1)\)
has \(\pi_1M=\mathbb Z/2\) and \(H_1(M;\mathbb Z)=\mathbb Z/2\).
It is not an integral homology sphere. Nevertheless every
finite-dimensional real or complex local system \(E\) satisfies
\[
 H^1(M;E)=H^1(\pi_1M;E)=0
\]
by averaging. The equality in degree one follows because the simply
connected universal cover gives an exact lifted cellular resolution
through the degrees needed for \(H^1\).
The ordinary homology calculation is also the cellular one in
\citet[Examples~2.42--2.43]{hatcher2002algebraic}.
\end{example}

\begin{proposition}[A homology sphere with a degenerate irreducible connection \ours]
There is a closed connected orientable integral homology
three-sphere \(M\) and an irreducible representation
\(\rho:\pi_1M\to\mathrm{SU}(2)\) such that
\(\dim_{\mathbb R}H^1(M;\operatorname{ad}\rho)=3\).
\end{proposition}
\begin{proof}
Let \(G=2I\subset\mathrm{SU}(2)\) be the binary icosahedral group
and \(Y=S^3/G\), with the free orientation-preserving action by
left multiplication on unit quaternions.
The order-\(120\) group has presentation
\(\langle a,b\mid a^5=b^3=(ab)^2\rangle\).
Its abelianization vanishes: taking the two relators as rows gives the integer
matrix \(\left(\begin{smallmatrix}5&-3\\-2&1\end{smallmatrix}\right)\),
of determinant \(-1\).
This presentation and calculation are recorded in
\citet[Example~2.38]{hatcher2002algebraic}.
Hence \(H_1(Y;\mathbb Z)=0\); orientability, connectedness and
Poincare duality give the integral homology of \(S^3\).

Set \(M=Y\#Y\). Van Kampen gives \(\pi_1M=G*G\);
abelianization and duality again give the integral homology of \(S^3\).
Define \(\rho\) to be the inclusion into \(\mathrm{SU}(2)\) on each
free factor. Its image is nonabelian and acts irreducibly on
\(\mathbb C^2\). For \(V=\mathfrak{su}_2\) with the adjoint action,
\(V^G=0\): the icosahedral rotation group has no common fixed axis.

A one-cocycle on \(G*G\) is uniquely determined by its restrictions
to the two free factors; extend by the cocycle identity along reduced
words. Averaging makes these restrictions
\(w_1-gw_1\) and \(w_2-gw_2\).
Since \(V^G=0\), the \(w_i\) are unique. Thus
\(Z^1(G*G;V)\cong V\oplus V\); principal cocycles form its diagonal
copy of \(V\). Consequently
\[
 H^1(M;\operatorname{ad}\rho)\cong
 (V\oplus V)/\operatorname{diag}V\cong V.
\]
The associated \(\mathrm{SU}(2)\) principal bundle is trivial, as
required in the source: over a three-dimensional CW complex a section
exists because \(\pi_0,\pi_1,\pi_2\) of \(\mathrm{SU}(2)\cong S^3\)
all vanish. This verifies the bundle and irreducibility conditions.
\end{proof}

\paragraph{Independent second check.}
Decompose \(M\) into two punctured copies of \(Y\) intersecting
in a collar of \(S^2\). With these coefficients each piece has
\(H^0=0\) and \(H^1=0\), while the intersection has \(H^0=V\)
and \(H^1=0\). Twisted Mayer--Vietoris gives
\(H^1(M;V)\cong V\) again. Puncturing does not change the
fundamental group or degree-one local-system cohomology.
This agrees with the cocycle quotient and computes the exact dimension;
it is stronger than merely exhibiting a family of representations.
